(1) 对态密度 \(g(\omega)=9N\omega^2/\omega_D^3\) 积分，\(\int_0^{\omega_D}g(\omega)d\omega=9N\omega_D^3/(3\omega_D^3)=3N\)，对应每个原子三个振动自由度。定义Debye温度 \(\Theta_D=\hbar\omega_D/k_B\)，将温度尺度与谱截止联系起来。
(2) 每个频率为 \(\omega\) 的声子模平均热能为 \(\hbar\omega/(e^{\hbar\omega/(k_BT)}-1)\)，所以 \(U=\int_0^{\omega_D}g(\omega)\hbar\omega/[e^{\hbar\omega/(k_BT)}-1]d\omega\)。令 \(x=\hbar\omega/(k_BT)\)，得到 \(U=9Nk_BT(T/\Theta_D)^3\int_0^{\Theta_D/T}x^3/(e^x-1)dx\)。
(3) 在 \(T\ll\Theta_D\)，积分上限可延至无穷，标准积分 \(\int_0^\infty x^3/(e^x-1)dx=\pi^4/15\)。因此 \(U\simeq(3\pi^4/5)Nk_BT^4/\Theta_D^3\)，求温度导数得 \(C_V\simeq(12\pi^4/5)Nk_B(T/\Theta_D)^3\)，呈明确的三次幂规律。
(4) 当 \(T\gg\Theta_D\)，每个简谐模平均能量趋近 \(k_BT\)，由总模数 \(3N\) 得 \(C_V\to3Nk_B\)，即Dulong–Petit定律。低温采用积分上限无穷的步骤在高温不成立，不能把 \(T^3\) 关系外推到高温。
